\documentclass[letterpaper,11pt]{article}

\usepackage{latexsym}
\usepackage[empty]{fullpage}
\usepackage{titlesec}
\usepackage{marvosym}
\usepackage[usenames,dvipsnames]{color}
\usepackage{verbatim}
\usepackage{enumitem}
\usepackage[hidelinks]{hyperref}
\usepackage{fancyhdr}
\usepackage[english]{babel}
\usepackage{tabularx}
\input{glyphtounicode}


%----------FONT OPTIONS----------
% sans-serif
% \usepackage[sfdefault]{FiraSans}
% \usepackage[sfdefault]{roboto}
% \usepackage[sfdefault]{noto-sans}
% \usepackage[default]{sourcesanspro}

% serif
% \usepackage{CormorantGaramond}
% \usepackage{charter}


\pagestyle{fancy}
\fancyhf{} % clear all header and footer fields
\fancyfoot{}
\renewcommand{\headrulewidth}{0pt}
\renewcommand{\footrulewidth}{0pt}

% Adjust margins
\addtolength{\oddsidemargin}{-0.5in}
\addtolength{\evensidemargin}{-0.5in}
\addtolength{\textwidth}{1in}
\addtolength{\topmargin}{-.5in}
\addtolength{\textheight}{1.0in}

\urlstyle{same}

\raggedbottom
\raggedright
\setlength{\tabcolsep}{0in}

% Sections formatting
\titleformat{\section}{
  \vspace{-10pt}\scshape\raggedright\large %%%%%%%%%%%%%%%%%%%%%%%%%%%%%%%%%%%%%%%%%%%%%%%%
}{}{0em}{}[\color{black}\titlerule \vspace{-5pt}]

% Ensure that generate pdf is machine readable/ATS parsable
\pdfgentounicode=1

%-------------------------
% Custom commands
\newcommand{\resumeItem}[1]{
  \item\small{
    {#1 \vspace{-2pt}}
  }
}

\newcommand{\resumeSubheading}[4]{
  \vspace{-2pt}\item
    \begin{tabular*}{0.97\textwidth}[t]{l@{\extracolsep{\fill}}r}
      \textbf{#1} & #2 \\
      \textit{\small#3} & \textit{\small #4} \\
    \end{tabular*}\vspace{-7pt}
}

\newcommand{\resumeSubSubheading}[2]{
    \item
    \begin{tabular*}{0.97\textwidth}{l@{\extracolsep{\fill}}r}
      \textit{\small#1} & \textit{\small #2} \\
    \end{tabular*}\vspace{-7pt}
}

\newcommand{\resumeProjectHeading}[2]{
    \item
    \begin{tabular*}{0.97\textwidth}{l@{\extracolsep{\fill}}r}
      \small#1 & #2 \\
    \end{tabular*}\vspace{-7pt}
}

\newcommand{\resumeSubItem}[1]{\resumeItem{#1}\vspace{-2pt}}

\renewcommand\labelitemii{$\vcenter{\hbox{\tiny$\bullet$}}$}

\newcommand{\resumeSubHeadingListStart}{\begin{itemize}[leftmargin=0.15in, label={}]}
\newcommand{\resumeSubHeadingListEnd}{\end{itemize}}
\newcommand{\resumeItemListStart}{\begin{itemize}}
\newcommand{\resumeItemListEnd}{\end{itemize}\vspace{-5pt}}

%-------------------------------------------
%%%%%%  RESUME STARTS HERE  %%%%%%%%%%%%%%%%%%%%%%%%%%%%


\begin{document}

\begin{center}
    \textbf{\Huge \scshape Andy Setiawan} \\ \vspace{1pt}
    \small +1 604-441-7964 $|$ \href{mailto:andy.setiawan9910@gmail.com}{\underline{andy.setiawan9910@gmail.com}} $|$ 
    \href{https://www.linkedin.com/in/andysetiawan1405/}{\underline{linkedin.com/in/andysetiawan1405}} $|$
    \href{https://fixables.github.io/}{\underline{fixables.github.io}}
\end{center}
\vspace{-10pt}
%-----------TECHNICAL SKILLS-----------
\section{Technical Skills}
 \begin{itemize}[leftmargin=0.15in, label={}]
    \small{\item{
     
     \textbf{Programming}{: C/C++, Python, MATLAB, ARM Assembly, Verilog, Arduino (IDE/PlatformIO)} \\
     
     \textbf{Embedded \& Firmware}{: ESP32, FreeRTOS, 8051, UART/SPI/I2C, PWM, Microbit, Circuit Playground, Scratch} \\
     
     \textbf{CAD \& Simulation}{: Altium Designer, LTSpice, Simulink, Quartus, ModelSim, OnShape, Tinkercad} \\
     
     \textbf{Instrumentation}{: Soldering (SMT \& THT), Oscilloscope, Multimeter, Function Generator, Power Supply} \\
     
     \textbf{Fabrication}{: Soldering (SMT \& THT), Drilling, Filing, Wiring, Mechanical Assembly} \\
     
     \textbf{Collaboration}{: Git, Microsoft Excel/Word/PowerPoint, Teams, Documentation \& Reporting, Project Management}
     }
    }
 \end{itemize}

%-----------EDUCATION-----------
\section{Education}
  \resumeSubHeadingListStart
    \resumeSubheading
      {The University of British Columbia}{Vancouver, BC}
      {Bachelor of Applied Science in Electrical Engineering}{Sept. 2023 -- May 2027}
      \resumeItemListStart
        \resumeItem{Indonesia Maju Scholarship Awardee (\textbf{CAD 400,000+ full-ride} from the MoECRT of Indonesia) $|$ Dean's List} 
        \resumeItem{Faculty of Applied Science International Student Scholarship $|$ Kenneth George VanSacker Memorial Scholarship}
      \resumeItemListEnd
  \resumeSubHeadingListEnd

%-----------EXPERIENCE-----------
\section{Technical Experience}
  \resumeSubHeadingListStart

    \resumeSubheading
      {Electrical Technician Apprentice \& Customer Service Representative}{Jan. 2016 -- Sept. 2023}
      {Tjahya Elektronik, Electronics Workshop \& Supplies}{Kintamani, Bali}
      \resumeItemListStart
        % \resumeItem{Performed \textbf{AC/DC measurements} (voltage, current, continuity, resistance) using multimeters and oscilloscopes to isolate failures such as regulator dropouts and shorted components, ensuring measurement accuracy.}
        
        \resumeItem{Followed structured troubleshooting workflows (visual inspection, staged power-up, signal probing, component substitution) and maintained clear service notes of symptoms, test results, and replaced parts}
        
        \resumeItem{Diagnosed and repaired 100+ \textbf{electrical devices}  including CRT and LED televisions, speakers, and household appliances using systematic debugging to \textbf{restore full functionality} and extend device lifespan by up to 50\%.}
        
        \resumeItem{Upgraded 20+ \textbf{legacy speaker systems} by integrating \textbf{Bluetooth modules}, \textbf{redesigning control PCBs}, and modifying casings, enhancing usability and extending product value.}
           
        % \resumeItem{Calibrated and installed digital and satellite TV receivers to detect \textbf{Indonesian communication satellites}.} 
        % %(Ku-Band, Telkom 3S, Palapa-D, Measat)

        % \resumeItem{Developed an all-in-one \textbf{USB configuration tool} for satellite TV receivers that \textbf{automated firmware installation} and channel setup, reducing configuration time by 80\%.}
        
        % \resumeItem{Managed store operations including \textbf{inventory tracking, component pricing, and supplier coordination}, maintaining stock of 200+ electronic components and 100+ unique household electrical appliances.}
        
        % \resumeItem{Provided \textbf{technical recommendations and on-site support} to customers, achieving a 90\% satisfaction rate through clear communication and personalized service.}
        
        % \resumeItem{Currently provide online consultation for \textbf{remote troubleshooting and component recommendations} to support returning customers and hobbyists.}
      \resumeItemListEnd

    \resumeSubheading
      {ELEC 201 Teaching Assistant}{Sep. 2025 -- Present}
      {UBC Department of Electrical and Computer Engineering}{Vancouver, BC}
      \resumeItemListStart
        \resumeItem{Led lab sessions for 50+ students in \textbf{circuit fundamentals}, SC/OC tests, Thevenin equivalents, circuit design, power calculation, and breadboard troubleshooting using power supply, oscilloscope, and signal generator.}
        \resumeItem{Mentored students on nonlinear and active components such as \textbf{transistors}, \textbf{diodes}, \textbf{rectifiers}, \textbf{555 timer oscillators}, and \textbf{operational amplifiers}, comparing real vs. theoretical circuit to build intuition.}
        \resumeItem{Invigilated midterm examinations and assisted with \textbf{WebWork} assessments, ensuring academic integrity, consistent grading standards, and a fair testing environment.}
    \resumeItemListEnd

    \resumeSubheading
      {Technician \& Facilitator}{May 2025 -- Sept. 2025}
      {Steamoji Kerrisdale}{Vancouver, BC}
      \resumeItemListStart
        \resumeItem{\textbf{Repaired and maintained} digital fabrication tools including \textbf{3D pens}, \textbf{3D printers}, and a \textbf{Glowforge laser cutting machine}, saving over \$1,000 in repair and replacement costs.}
        \resumeItem{Taught and mentored over 50 students (ages 5-14) during summer camps and regular sessions in \textbf{Arduino, robotics, programming, soldering, 3D design, electronics circuit, and digital fabrication}.}
      \resumeItemListEnd

  \resumeSubHeadingListEnd

%-----------Design Teams-----------
\section{Engineering Design Teams}
  \resumeSubHeadingListStart
    \resumeSubheading
      {Open Robotics, Firmware Co-Lead}{Sept. 2023 -- Present}{}{}
       \vspace{-20pt}
       

% %%%%%%%%% 4 Bullet points version
%         \resumeProjectHeading
%         {\textbf{} \emph{Altium, LTSpice, Control System, Motor Control, RTOS, Power Distribution, Hardware Validation, Load Calculation}}{}
%        \resumeItemListStart
%           \resumeItem{Built a \textbf{haptic BLDC knob system} that lets you physically feel what is happening in basic electrical elements and circuits (R, L, C, RLC, and a diode) by turning a motorized knob with encoder feedback, and a small UI, translating high-level interaction goals into a real embedded control flow.}
          
%           \resumeItem{Developed \textbf{ESP32} firmware using \textbf{dual-core FreeRTOS} to coordinate motor control, sensor sampling, and SPI OLED display 
%           tasks with deterministic timing and clean task separation.}
          
%           \resumeItem{Designed a \textbf{robot emergency-stop power cutoff} system for Robocup@Home, creating a safety architecture that isolates power quickly during fault conditions.}
          
%           \resumeItem{Completed \textbf{schematic capture and 2-layer PCB layout} in \textbf{Altium}, then validated cutoff timing with \textbf{LTSpice} and custom test plans, achieving less than 50 ms from button actuation to full isolation.}
%         \resumeItemListEnd


%%%%%%%%%%%%%%%%%% 2 projects separated version
        \resumeProjectHeading
          {\textbf{Haptic Knob} $|$ \emph{C++, C, Control System, PID, FOC, ESP32, Motor Control, RTOS}}{Ongoing Project}
          
        \resumeItemListStart
        % \resumeItem{Contributed to \textbf{high-level system architecture} for both hardware and firmware, defining control flow, sensing strategy, and motor–UI integration for a haptic BLDC system.}
        
        \resumeItem{Developed \textbf{ESP32} firmware using \textbf{dual-core FreeRTOS} to coordinate motor control, sensor sampling, and SPI OLED display 
        tasks with deterministic timing and clean task separation.}
        
        \resumeItem{Implemented closed-loop \textbf{BLDC motor control} using \textbf{current sensing}, \textbf{encoder feedback}, and \textbf{PID and Field  Oriented Control algorithms} to generate haptic torque that emulates different circuit behaviors.}

        \resumeItem{Generated structured \textbf{technical documentation} including system block diagram, I/O definitions, signal mapping, control flow descriptions, and \textbf{test procedures} to support cross-team integration and system validation.}

        \resumeItemListEnd
% \vspace{-5pt}
%         \resumeProjectHeading
%         {\textbf{Robocup@Home} $|$ \emph{Altium, LTSpice, Power Distribution, Hardware Validation, Load Calculation}}{}

%         \resumeItemListStart
%         \resumeItem{Designed the robot \textbf{power safety architecture} and completed \textbf{schematic capture + 2-layer PCB layout} in \textbf{Altium Designer} with relays and protection circuitry for immediate \textbf{fail-safe power isolation}.}
        
%         \resumeItem{Validated cutoff performance using \textbf{LTSpice simulations} and \textbf{custom testbenches}, achieving less than 50 ms response time from emergency-button actuation to full power isolation.}

        % \resumeItem{Integrated safety circuitry with \textbf{motor controllers} and \textbf{communication modules} in collaboration with \textbf{electrical, software, drivetrain, and gripper subteams.}}
      \resumeItemListEnd
      
    % \resumeSubheading
    %   {Biological Internet of Things, Instrumentation Team}{Sept. 2025 -- Present}{}{}
    %   \vspace{-20pt}
    
    %   \resumeProjectHeading
    %     {\textbf{} \emph{Altium, Voltage \& Data Line Isolation, Sensor Integration}}{}
    %   \resumeItemListStart
    %     \resumeItem{Designed a \textbf{2-layer motherboard PCB} in \textbf{Altium Designer}, using a modular architecture with an \textbf{ESP32 carrier} and plug-in sensor boards for temperature, dissolved oxygen, pressure, and pH.}
    %     \resumeItem{Implemented voltage and data-line isolation using the \textbf{ADM3260} (\textbf{isolated I2C} with \textbf{integrated isolated power}) to protect the MCU and improve signal robustness across sensor modules.}
    %   \resumeItemListEnd

  % \resumeSubHeadingListEnd

%-----------PROJECTS-----------
\section{Projects}
  \resumeSubHeadingListStart
    \resumeProjectHeading
      {\textbf{RoboMaestro} $|$ \emph{H-Bridge, Gate Drivers, PWM, Motor Control, Power Electronics, Circuit Design}}{Ongoing}
      \resumeItemListStart
        \resumeItem{Designed a \textbf{discrete MOSFET H-bridge} motor driver and \textbf{2-layer PCB in Altium}, focusing on gate-drive integrity, level shifting, and noise-aware layout for reliable \textbf{PWM operation} and \textbf{dead-time control}.}
        
        \resumeItem{Capturing and routing the \textbf{2 layer motor driver PCB in Altium}, implementing \textbf{gate-drive networks with decoupling suppression} to reduce noise and improve switching robustness.}
      \resumeItemListEnd

    % \resumeProjectHeading
    %   {\textbf{Smart Energy Meter} $|$ \emph{ESP32, IoT, Power Measurement, Sensors, Python Analytics}}{Ongoing Project}
    %   \resumeItemListStart
    %     \resumeItem{Designing an \textbf{IoT-enabled energy monitoring system} using \textbf{ESP32}, SCT-013 current sensor, and ZMPT101B voltage module to compute real-time RMS voltage/current, power, and power factor with a TFT display.}
    %     % \resumeItem{Developing robust data acquisition and calibration logic to improve measurement repeatability across loads and reduce noise sensitivity in live readings.}
    %     % \resumeItem{Collaborating with a Statistics student to build a \textbf{Python analytics pipeline} for visualization, anomaly detection, and short-term forecasting using \textbf{ARIMA} with interactive dashboards.}
    %   \resumeItemListEnd
      
    % \resumeProjectHeading
    %   {\textbf{Coin Picking Robot} $|$ \emph{EFM8LB1, MSP430, LPC824, PWM, Wireless, Sensors, Embedded Systems}}{Mar. 2025}
    %   \resumeItemListStart
    %     \resumeItem{Engineered a \textbf{multi-microcontroller robotic system} with \textbf{JDY-40 wireless} communication, supporting autonomous coin collection and manual remote control via joystick with LCD feedback.}
        
    %     \resumeItem{Designed and implemented the robot’s \textbf{control logic and finite state machine (FSM)} to execute coin-picking behavior, integrating sensor inputs with actuator outputs using \textbf{PWM motor control} and \textbf{ISRs} for real-time response.}

    %     % \resumeItem{Developed the \textbf{motor control firmware} including \textbf{PWM drive}, \textbf{interrupt service routines (ISRs)} for real-time sensing/control, and an \textbf{autonomous state machine} coordinating navigation and coin-pickup behavior.}

    %     % \resumeItem{Integrated robustness features including an \textbf{ultrasonic obstacle detector}, light-sensitive headlamp, laser alignment guide, coin-counting laser, and turn-signal LEDs for improved usability and reliability.}
    %   \resumeItemListEnd

    % \resumeProjectHeading
    %   {\textbf{Reflow Oven Controller} $|$ \emph{8051, FSM, PWM, Thermocouple, Op-Amp, Python, Calibration}}{(Feb. 2025)}
    %   \resumeItemListStart
    %     % \resumeItem{Programmed an \textbf{8051-based FSM} reflow controller with \textbf{PWM-driven SSR} switching, achieving \textbf{$\pm$2\,$^\circ$C} thermal accuracy using a K-type thermocouple front-end with op-amp amplification}
    %     % \resumeItem{Designed a user-friendly interface with LCD + buttons, adjustable reflow parameters, safety LEDs, and an automatic abort feature to support reliable SMD soldering workflows.}
    %     % \resumeItem{Added usability enhancements including a completion jingle and animated LCD messages to improve operator feedback and reduce setup errors during reflow runs.}
    %     % \resumeItem{Designed a \textbf{K-type thermocouple signal-conditioning circuit} (op-amp amplification + filtering) and implemented \textbf{8051} temperature sampling for reliable sensor readings.}
    %     \resumeItem{Designed a \textbf{K-type thermocouple signal-conditioning circuit} (op-amp amplification/filtering) and built \textbf{Python/Excel logging} with live \textbf{strip-chart visualization} and calibration, achieving ±2°C temperature accuracy}
    %   \resumeItemListEnd

    % \resumeProjectHeading
    %   {\textbf{Oscilloscope \& Multimeter} $|$ \emph{555 Timer, C, Embedded Measurement, Python Plotting}}{(Jan. 2025)}
    %   \resumeItemListStart
    %     \resumeItem{Designed a \textbf{555 timer-based measurement circuit} paired with a microcontroller to probe voltage, resistance, and capacitance via frequency-shift sensing, achieving approximately \textbf{1 mV}, \textbf{10 $\Omega$}, and \textbf{1 pF} resolution.}
    %     \resumeItem{Programmed in \textbf{C} to capture oscillator frequency, compute electrical parameters, and handle measurement modes with stable sampling and filtering.}
    %     \resumeItem{Built a \textbf{Python} visualization workflow for real-time plotting to verify measurement stability and support quick calibration during testing.}
    %   \resumeItemListEnd

  % \resumeSubHeadingListEnd



%


%-------------------------------------------
\end{document}
